\chapter{Teoretične osnove in sorodna dela}
\label{ch:teorija}

Poglavje predstavi osnove, potrebne za razumevanje celotne poti od
zvoka v prostoru do prikazane odločitve. Najprej je opisan prehod iz zvočnega
valovanja v zaporedje števil, nato pretvorba tega zaporedja v časovno-frekvenčno
predstavitev. Sledijo razvrščanje zvočnih dogodkov, zasnova modela YAMNet,
prenosno učenje, celoštevilska kvantizacija in izvajanje na robni napravi.
Poglavje se zaključi s primerjavo sorodnih pristopov.

\section{Digitalni zvok in zajem signala}

Zvok je sprememba zračnega tlaka, ki se širi po prostoru. Mikrofon to pretvori
v električni signal, ki se za obdelavo na mikrokrmilniku vzorči in predstavi kot
zaporedje diskretnih vrednosti. Posamezna vrednost je \emph{vzorec}, število
vzorcev v eni sekundi pa je \emph{frekvenca vzorčenja}.

V izvedenem sistemu je frekvenca vzorčenja \SI{16}{\kilo\hertz}, zato je po Nyquistovem
pravilu mogoče brez prekrivanja predstaviti frekvence do polovice frekvence vzorčenja, oziroma
do \SI{8}{\kilo\hertz}. To zadošča uporabljeni predobdelavi, ki upošteva območje
od \SI{125}{\hertz} do \SI{7.5}{\kilo\hertz}.

Posamezen vzorec je trenutna številska vrednost digitalnega signala. Sistem
uporablja predznačene šestnajstbitne vzorce z ničlo kot referenčno vrednostjo;
dovoljeno celoštevilsko območje je od $-32.768$ do $32.767$. Večja absolutna
vrednost oziroma amplituda praviloma pomeni močnejši trenutni signal, vendar sama vrednost vzorca
še ne pove, za kateri zvok gre.

Na razvojni plošči je digitalni mikroelektromehanski mikrofon (MEMS). Njegov neposredni
izhod ni šestnajstbitna amplituda, temveč hiter enobitni pulzno-gostotni tok.
Razmerje med enicami in ničlami predstavlja amplitudo zvočnega signala.
Digitalni filter enobitni tok pretvori v običajne večbitne zvočne vzorce z
nižjo frekvenco vzorčenja, ki jih nato uporabi nadaljna predobdelava.
Nastavitve mikrofona, filtra in prenosa v pomnilnik so opisane v
poglavju~\ref{ch:implementacija}.

\section{Časovno-frekvenčna predstavitev}

Časovni potek zvočnega signala prikazuje spreminjanje amplitude, ne pokaže pa
neposredno, katere frekvence so prisotne v posameznem trenutku. Podatek o
frekvenčni vsebini pa je zelo pomemben za razpoznavanje okoljskih zvokov. Sirena tako
navadno vsebuje dlje trajajoče tone v ozkih frekvenčnih območjih. Razbitje
stekla je kratko in zajema širok razpon frekvenc, govor pa se spreminja skozi
čas. Zvočni signal zato pred obdelavo z nevronskim modelom pretvorimo v
strnjeni spektrogram, ki prikazuje spreminjanje njegove frekvenčne vsebine
skozi čas.

\subsection{Kratkočasovna Fourierjeva transformacija}

Fourierjeva transformacija pokaže, katere frekvence so prisotne v signalu. Če
bi jo izračunali za celoten posnetek, bi izgubili podatek, kdaj so se posamezne
frekvence pojavile. Pri kratkočasovni Fourierjevi transformaciji se zato signal razdeli na kratka,
delno prekrivajoča se okna. Transformacija se izračuna za vsako okno posebej.

Za časovno okno z zaporedno številko $m$ in frekvenčno mesto $k$ je izračun
zapisljiv kot

\begin{equation}
  X(m,k) = \sum_{n=0}^{N-1} x[n+mH]\,w[n]\,
  \mathrm{e}^{-\mathrm{j}2\pi kn/N_F},
  \label{eq:stft}
\end{equation}

kjer je $x$ vhodni signal, $w$ utežna okenska funkcija, $N$ število dejanskih
vzorcev v oknu, $H$ korak med začetkoma dveh oken in $N_F$ velikost
Fourierjeve transformacije. Kompleksni rezultat $X(m,k)$ vsebuje velikost in
fazo posamezne frekvenčne sestavine. Za nadaljnjo predobdelavo je pomembna
predvsem velikost oziroma amplituda.

Izbrana različica modela YAMNet določa dolžino in razmik časovnih oken. Vsako
okno vsebuje 400 vzorcev oziroma \SI{25}{\milli\second} zvoka, začetek
naslednjega pa je zamaknjen za 160 vzorcev oziroma
\SI{10}{\milli\second}. Sosednji okni si tako delita 240 vzorcev oziroma
\SI{15}{\milli\second}. Prekrivanje zmanjša vpliv položaja kratkega dogodka glede
na meje oken.

Vsako okno najprej pomnožimo s Hannovo okensko funkcijo, ki ublaži ostre
prehode na robovih okna. Okno s 400 vzorci nato dopolnimo s 112 ničlami do
dolžine 512 vzorcev. Dopolnitev ne ustvari novih informacij, omogoči pa učinkovit
izračun Fourierjeve transformacije in zgosti mrežo frekvenčnih mest.

En vhod modela vsebuje 96 zaporednih časovnih oken. Prvo okno potrebuje 400
vzorcev, vsako naslednje pa se začne 160 vzorcev pozneje. Celoten potreben
odsek zato vsebuje

\begin{equation}
  400 + (96-1)\cdot160 = 15.600
\end{equation}

vzorcev oziroma \SI{975}{\milli\second} zvoka. Dokumentacija modela časovno
dolžino navadno podaja kot $96\cdot\SI{10}{\milli\second}=
\SI{960}{\milli\second}$, torej kot število oken, pomnoženo z njihovim korakom.
Dejanska izvedba mora upoštevati še trajanje prvega okna, zato za izračun vseh
96 oken potrebuje 15.600 vzorcev. Nova odločitev se v končni aplikaciji izračuna vsakih
\SI{480}{\milli\second}, saj si zaporedna vhoda delita približno polovico
zvočnega gradiva.

\subsection{Melodična lestvica in strnjeni spektrogram}

Frekvenčna mesta Fourierjeve transformacije so enakomerno razporejena v
hercih. Človeško zaznavanje višine pa ni linearno: pri nizkih frekvencah so
majhne razlike zaznavnejše kot pri visokih. Melodična lestvica frekvenco približno
preslika v zaznavno primernejšo razdaljo. Ena izmed običajnih preslikav je

\begin{equation}
  M(f) = 2595\log_{10}\left(1+\frac{f}{700}\right),
  \label{eq:mel}
\end{equation}

kjer je $f$ frekvenca v hercih, $M(f)$ pa položaj na melodični lestvici.

Za vsako časovno okno Fourierjeva transformacija izračuna amplitude v
posameznih frekvenčnih točkah. Tako dobljeni spektri zaporednih oken tvorijo
spektrogram. Na spekter vsakega okna nato uporabimo 64 delno prekrivajočih se
trikotnih filtrov. Vsak filter združi vrednosti več sosednjih frekvenčnih točk
v eno vrednost, ki predstavlja amplitudo v določenem melodičnem pasu. Za vsako
časovno okno tako dobimo vektor 64 značilk.

Na značilke nato uporabimo logaritem, ker človeški sluh jakost zvoka zaznava
približno logaritemsko glede na amplitudo signala. Logaritem hkrati zmanjša
razpon med šibkimi in močnimi komponentami. Zaporedje 96 logaritmiranih
vektorjev tvori matriko velikosti $64\times96$, ki jo v tej nalogi imenujemo
\emph{strnjeni spektrogram}. Ta matrika je vhod v nevronski model. Enake
parametre uporablja uradna različica YAMNet-1024 za STM32N6
\cite{ststm32aimodelzooyamnet}.

\section{Razvrščanje zvočnih dogodkov}

Model prejme zaporedje 96 vektorjev značilk in izračuna oceno za vsak naučeni
razred. Med učenjem je za vsak primer znana pravilna oznaka razreda, postopek
učenja pa uteži modela prilagaja tako, da je ocena pravilnega razreda čim večja.
Ob izvajanju model za nov vhod vrne vektor ocen. Razred z največjo oceno je
trenutna najverjetnejša razlaga vhodnega odseka, ni pa nujno dovolj zanesljiv za
prikaz opozorila.

Tudi kadar so izhodi normalizirani in je njihova vsota ena, vrednosti niso
samodejno umerjene verjetnosti pravilnosti. Ocena 0,90 pomeni, da model močno
daje prednost določenemu razredu glede na druge naučene razrede. Ne pomeni pa,
da bo napoved pravilna v natanko 90 odstotkih enakih primerov. Zato so v nalogi
vrednosti poimenovane \emph{ocene modela}; pravilnost sistema se meri ločeno na
preskusnih posnetkih.

V resničnem okolju se lahko pojavijo tudi zvoki zunaj nabora razredov,
uporabljenih pri učenju modela. Med uporabo lahko model prejme glasbo, trkanje,
premikanje predmetov ali drug zvok, ki ne spada v noben ciljni razred. Če ima
model samo nevarnostne izhode, mora tudi
takemu vhodu dodeliti enega izmed njih. Razreda \emph{govor} in \emph{ostalo} zato delujeta kot
varovalna razreda. \emph{Govor} pokriva zelo pogost nenevaren zvok, razred \emph{ostalo} pa
raznolike preostale primere. Razred \emph{ostalo} kljub temu ne more predstavljati
vseh mogočih neznanih zvokov; njegova učinkovitost je odvisna od raznolikosti
učnih podatkov.

Posamezna napoved modela prav tako še ni končna odločitev naprave. Ocene se
med zaporednimi, prekrivajočimi se odseki lahko spreminjajo. Odločitvena plast
zato uporablja glajenje, razredne pragove in zahtevano število zaporednih
potrditev. Če noben razred ne doseže svojega pogoja, sistem zvok zavrne kot
\emph{negotovo napoved}. Višji prag zmanjša število lažnih alarmov, vendar lahko poveča število
spregledanih dogodkov. Nižji prag ima obraten učinek. Končni sistem je zato
treba vrednotiti kot celoto: zajem, predobdelavo, model in časovno odločanje,
ne samo kot nevronsko mrežo na datotekah.

\section{Prilagoditev in kvantizacija modela YAMNet}

Razdelek najprej predstavi prednaučeni model YAMNet in njegov izhod, nato pa
pojasni, kako smo model s prenosnim učenjem prilagodili izbranim razredom in ga z
osembitno kvantizacijo pripravili za izvajanje na mikrokrmilniku.

\subsection{Prednaučen model YAMNet}

YAMNet je prednaučeni model, ki vhodni zvok razvršča med 521 razredov zbirke
AudioSet \cite{hershey2017cnn,tensorflowyamnetrepo}, ki je obsežna,
hierarhično označena zbirka zvočnih dogodkov iz spletnih videoposnetkov
\cite{gemmeke2017audioset}. Model zato že prepoznava splošne vzorce govora,
živali, glasbe, strojev in drugih okoljskih zvokov.

YAMNet temelji na arhitekturi MobileNetV1 \cite{howard2017mobilenets}. Njena
ključna lastnost so globinsko ločljive konvolucije (angl. \emph{depthwise
separable convolutions}). Običajna konvolucija v enem koraku išče prostorske
vzorce in hkrati povezuje vse vhodne kanale. Globinsko ločljiva konvolucija
nalogo razdeli na dva dela. Najprej globinska konvolucija (angl.
\emph{depthwise convolution}) uporabi ločen filter za vsak kanal. Nato
enotočkovna konvolucija (angl. \emph{pointwise convolution}) velikosti
$1\times1$ poveže rezultate med kanali. Tak razcep praviloma potrebuje manj
parametrov in operacij, zato je primeren za mobilne in vgrajene naprave.

\subsection{Prenosno učenje}

Pri prenosnem učenju (angl. \emph{transfer learning}) prednaučeni model prilagodimo 
za reševanje nove, specifične naloge. Namesto da bi nevronsko mrežo učili povsem od 
začetka, obdržimo arhitekturo in že naučene uteži večjega dela modela, ki služi kot 
ekstraktor značilk. Zamenjamo in na novo naučimo le zadnjo, polno povezano izhodno 
plast. Na ta način model obdrži zmožnost prepoznavanja splošnih vzorcev (npr. zvočnih 
značilk v modelu YAMNet), nove parametre pa prilagodimo izbranim razredom. Ker se pri 
tem posodablja le majhen del parametrov, postopek bistveno zmanjša potrebo po obsežnih 
učnih množicah in skrajša čas učenja. Podrobnosti o tem, kako smo ta postopek uporabili 
v naši nalogi, so opisane v poglavju~\ref{ch:implementacija}.

\subsection{Osembitna celoštevilska kvantizacija}

Model se med učenjem praviloma računa z 32-bitnimi števili s plavajočo vejico.
Na mikrokrmilniku tak zapis poveča velikost uteži in ni najprimernejši za
namenski nevronski pospeševalnik. Pri kvantizaciji se realna vrednost $r$
približno preslika v osembitno celo število $q$ z zvezo

\begin{equation}
  q = \operatorname{clip}\!\left(\operatorname{round}\!\left(\frac{r}{s}\right)
      + z\right),
  \label{eq:kvantizacija}
\end{equation}

kjer je $s$ pozitivni kvantizacijski korak, $z$ pa ničelna točka. Funkcija
$\operatorname{clip}$ rezultat omeji na dovoljeno osembitno območje. Korak
določa, kolikšen razpon realnih vrednosti predstavlja ena celoštevilska
stopnja, ničelna točka pa poskrbi, da je mogoče natančno predstaviti realno
ničlo. Pri povratni pretvorbi velja približek $r\approx s(q-z)$.

Ker ima osembitna utež štirikrat manjši zapis od 32-bitne uteži, se bistveno
zmanjša pomnilniški odtis uteži. Nevronski pospeševalnik je načrtovan
prav za izvajanje takšnih celoštevilskih operacij. Vendar pa kvantizacija vrednosti
zaokrožuje, kar lahko zmanjša kakovost modela. Izgubo kakovosti zaradi napake zaokroževanja 
omejimo s kalibracijo: na manjšem vzorcu učnih podatkov izmerimo dinamični razpon vmesnih
aktivacij, na katerega nato optimalno preslikamo 8-bitno lestvico. Razmerje med točnostjo, 
hitrostjo in velikostjo pa je treba preveriti z meritvami, ne zgolj napovedati \cite{jacob2018quantization}.

\section{Robna umetna inteligenca in nevronski pospeševalniki}

Pri obdelavi na robni napravi (angl. edge) se podatki analizirajo neposredno pri svojem izvoru.
V obravnavanem sistemu celotna veriga od zajema zvoka, predobdelave in nevronskega sklepanja do 
prikaza na zaslonu poteka izključno na lokalni razvojni plošči. Surovi zvočni vzorci se obdelajo
sproti v delovnem pomnilniku in se takoj po analizi prepišejo z novimi podatki. Takšna izvedba 
prinaša dve ključni prednosti: zagotavlja visoko stopnjo zasebnosti uporabnikov, saj posnetki 
nikoli ne zapustijo naprave, ter omogoča hiter odziv in zanesljivo delovanje neodvisno od omrežne povezave.

Procesorsko jedro Arm Cortex-M55 je primerno za krmiljenje
periferije, izračun strnjenega spektrograma, odločanje ter uporabniški vmesnik. Izvajanje
več milijonov ponavljajočih se konvolucijskih operacij samo na njem pa ni
najučinkovitejša uporaba strojne opreme. Nevronski pospeševalnik Neural-ART v
mikrokrmilniku STM32N6 je namenjen vzporednemu izvajanju takšnih operacij.

\clearpage
\section{Sorodna dela}

Teoretičnemu ozadju sledi pregled sorodnih raziskav, ki omogoča neposredno primerjavo 
in umestitev našega sistema med obstoječe pristope.

\subsection{Sorodni primeri za naprave z omejenimi viri}

Mohaimenuzzaman in sodelavci so razvili model Micro-ACDNet in postopek, s
katerim so večjo konvolucijsko mrežo za razpoznavanje surovega zvoka stisnili
in kvantizirali \cite{mohaimenuzzaman2023edge}. Njihova različica za 50 razredov ima približno
131.000 parametrov in je bila izvedena na mikrokrmilniški plošči Sony Spresense.
Študija poudari omejitve delovnega pomnilnika in pokaže, da lahko kvantizacija
opazno spremeni pravilnost. V primerjavi s to nalogo je njihov model manjši in vhod
obdeluje neposredno kot zvočno obliko, cilj pa je splošen postopek stiskanja,
ne namenski sistem vidnega opozarjanja.

Elliott in sodelavci so raziskali majhne modele za
razpoznavanje okoljskih zvokov \cite{elliott2021tiny}. Najmanjši model za šest
pisarniških razredov ima 6.642 parametrov in je pri preverjanju na telefonu
Samsung Galaxy S9 sklepanje opravil v približno \SI{7}{\milli\second}. Avtorji
so velikost izbrali z mislijo na pomnilnik mikrokrmilniške plošče Arduino Nano
33 BLE Sense, vendar fizične izvedbe na tej plošči niso predstavili. Delo
pokaže, da velikost modela ni edina možna pot: na izbranem ozkem problemu lahko
zelo majhen model tekmuje z večjim konvolucijskim modelom. Hkrati rezultati niso
neposredno primerljivi s to nalogo, saj se razlikujejo razredi, dolžina vhoda,
strojna platforma in preskusni podatki.

\subsection{Pomožni sistemi za gluhe in naglušne osebe}

Jain in sodelavci so razvili domači sistem HomeSound, sestavljen iz več
zaslonov z mikrofoni in pozneje tudi pametne ure. Prva različica je prikazovala
glasnost, višino, trajanje in prostor izvora zvoka, druga pa je dodala samodejno
razvrščanje 19 gospodinjskih zvokov. Obe različici so po tri tedne preizkušali
v domovih gluhih in naglušnih oseb; skupaj je sodelovalo šest različnih
gospodinjstev \cite{jain2020homesound}. Sistem je povečal zavedanje o dogajanju
v domu, vendar so uporabnike motile napačne razvrstitve in prepogoste vibracije
pametne ure. Pojavila so se tudi vprašanja zasebnosti obiskovalcev. Delo je za
to nalogo pomembno, ker pokaže, da visoka laboratorijska pravilnost sama ne
zagotavlja uporabnega opozorilnega sistema. Naša rešitev zato uporablja
časovno potrjevanje in razredne pragove.

Lin in sodelavci so predstavili dvostopenjski notranji sistem za prikaz zvokov
gluhim in naglušnim osebam \cite{lin2026indoor}. Manjši model Dynamic
MobileNet na računalniku Raspberry Pi najprej razvrsti zvok in oceni negotovost.
Če je ocena premalo zanesljiva, se posnetek pošlje strežniku, kjer ga obdela
večji model YAMNet. Rezultat se prikaže na pametnem telefonu ali uri. Sistem je
na preskusnem naboru, sestavljenem iz zbirk ESC-50 in ReaLISED, dosegel
90-odstotno točnost za 14 pomembnih razredov; za 86 širših skupin iz zbirke
AudioSet je dosegel 74 odstotkov. Rezultatov zaradi drugačnih razredov in
preskusnih podatkov ni mogoče neposredno primerjati z rezultati te naloge.
Arhitekturna razlika pa je jasna: dvostopenjski pristop s strežnikom omogoča
uporabo večjega modela ob negotovih primerih, medtem ko naša naprava celotno
obdelavo izvede lokalno in zvoka ne pošilja v omrežje.

\subsection{Primerjava in umestitev naloge}

\begin{table}[htbp]
  \centering
  \caption{Primerjava izbranih pristopov za razpoznavanje in prikaz okoljskih zvokov.}
  \label{tab:sorodna-dela}
  \scriptsize
  \setlength{\tabcolsep}{5pt}
  \renewcommand{\arraystretch}{1.3}
  \begin{tabular}{@{}>{\raggedright\arraybackslash}p{0.18\textwidth}
                   >{\raggedright\arraybackslash}p{0.18\textwidth}
                   >{\raggedright\arraybackslash}p{0.26\textwidth}
                   >{\raggedright\arraybackslash}p{0.28\textwidth}@{}}
    \toprule
    Pristop in vir & Platforma & Vhod in model & Poudarek \\
    \midrule
    Google YAMNet \cite{tensorflowyamnetrepo}
      & splošne računalniške platforme 
      & strnjeni spektrogram, MobileNetV1, 521 oznak 
      & splošno prednaučeno razpoznavanje zbirke AudioSet \\
    \addlinespace
    
    Micro-ACDNet \cite{mohaimenuzzaman2023edge}
      & Sony Spresense 
      & \SI{1.5}{\second} surovega zvoka, 50 razredov 
      & stiskanje, kvantizacija, pomnilnik in pravilnost \\
    \addlinespace
    
    Majhen transformator \cite{elliott2021tiny}
      & Samsung Galaxy S9 
      & melodični spektrogram, šest pisarniških razredov 
      & zelo majhen model; fizična izvedba na mikrokrmilniku ni bila prikazana \\
    \addlinespace
    
    HomeSound \cite{jain2020homesound}
      & domači zasloni in pametna ura 
      & značilnosti zvoka in pozneje 19 gospodinjskih razredov 
      & terenski preizkus; uporabnost, zasebnost in prepogosta opozorila \\
    \addlinespace
    
    Dvostopenjski sistem \cite{lin2026indoor}
      & Raspberry Pi in strežnik 
      & Dynamic MobileNet lokalno, YAMNet ob negotovih primerih 
      & širok nabor zvokov, vendar prenos zvoka do strežnika \\
    \addlinespace
    
    Ta naloga 
      & STM32N6570-DK in Neural-ART 
      & $64\times96$, YAMNet-1024, šest uporabniških razredov 
      & lokalni fizični sistem, nevarnosti, lažni alarmi, čas in pomnilnik \\
    \bottomrule
  \end{tabular}
\end{table}

Vrzel, ki jo naslavlja naloga, zato ni nov splošni model za razpoznavanje zvoka.
Prispevek je celovita prilagoditev prednaučenega modela in razvoj fizičnega
sistema za izbrane nevarne dogodke: neposredni zajem z vgrajenim mikrofonom,
enaka predobdelava pri učenju in na napravi, namenski učni razredi, lokalno
sklepanje, časovno stabilizirana odločitev, vidno opozorilo ter meritev na
končni strojni opremi. Sistem je zasnovan kot prototip pomoči gluhim in
naglušnim osebam, pri čemer običajno delovanje ne potrebuje računalniškega
oblaka in ne hrani surovega zvoka.